\documentclass[11pt, a4paper]{article}

% --- Paquetes Básicos ---
\usepackage[utf8]{inputenc}
\usepackage[spanish]{babel}
\usepackage{graphicx}
\usepackage{amsmath}
\usepackage{amsfonts}
\usepackage{hyperref}
\usepackage{geometry}
\usepackage{booktabs} % Para tablas prolijas
\usepackage{listings} % Para snippets de código si fuera necesario
\usepackage{xcolor}

\geometry{top=2.5cm, bottom=2.5cm, left=2.5cm, right=2.5cm}

% --- Configuración de Enlaces ---
\hypersetup{
    colorlinks=true,
    linkcolor=blue,
    filecolor=magenta,      
    urlcolor=cyan,
    pdftitle={TP2 SDN-NAT - Informe},
}

% --- Configuración para Código (Opcional) ---
\lstset{
    backgroundcolor=\color{gray!10},
    basicstyle=\ttfamily\small,
    breaklines=true,
    keywordstyle=\color{blue},
    commentstyle=\color{green!50!black},
    stringstyle=\color{red},
    frame=single,
    numbers=left,
    numberstyle=\tiny\color{gray},
    captionpos=b
}

% --- PORTADA ---
\title{
    \normalsize Universidad de Buenos Aires \\
    \normalsize Facultad de Ingeniería \\
    \normalsize Redes (TA048) / Introducción a los Sistemas Distribuidos (75.43) \\ [2cm]
    \Large \textbf{Trabajo Práctico N$^{\circ}$ 2: SDN - NAT} \\
    \large Implementación de NAT saliente con traducción por puertos (PAT) \\ [1cm]
}

\author{
    \textbf{Integrantes:} \\
    Apellido, Nombre - Padrón \\
    Apellido, Nombre - Padrón \\
    [2cm]
}
\date{23 de Junio de 2026} % Fecha de entrega pautada

\begin{document}

\maketitle
\thispagestyle{empty}
\clearpage
\tableofcontents
\clearpage

% --- SECCIÓN 1 ---
\section{Descripción General de la Solución}
\label{sec:descripcion_general}

\begin{small}
\color{blue}
\textbf{Nota Provisoria para el Desarrollo:} 
Acá deben explicar el enfoque macro del TP. Mencionen que se migró del modelo básico provisto (\texttt{protorouter.py}) donde solo se ruteaba estáticamente reescribiendo MACs, hacia un esquema de NAT con traducción de puertos (PAT). Expliquen cómo interactúan Mininet (plano de datos) y el controlador POX (plano de control) mediante OpenFlow para lograr que múltiples hosts de la red privada compartan una única IP pública (200.0.0.254).
\end{small}

\vspace{0.5cm}
[Escribir aquí la descripción general...]


% --- SECCIÓN 2 ---
\section{Extensión de la Topología}
\label{sec:extension_topologia}

Para satisfacer el escenario de múltiples clientes concurrentes en la red privada y evaluar correctamente el funcionamiento del NAT, se modificó el script original \texttt{topo.py}. En lugar de definir de manera estática una cantidad fija de hosts en la red privada, se parametrizó la creación de la topología para permitir una cantidad variable de los mismos.

Mediante el argumento de ejecución \texttt{-ph} (o \texttt{--private-hosts}), el usuario puede indicar la cantidad de hosts privados deseados al levantar la topología. El script se encarga de instanciar dinámicamente dichos hosts, asignándoles direcciones IP correlativas dentro de la subred privada (comenzando por la \texttt{192.168.1.2}) y direcciones MAC acordes a su número de host para simplificar su identificación (por ejemplo, \texttt{00:00:00:00:00:02} para el host con IP terminada en \texttt{.2}). Todos estos hosts se enlazan al mismo switch central \texttt{s1}.

Además, una modificación fundamental introducida en esta nueva topología es la eliminación completa de las configuraciones de ARP estático (anteriormente configuradas mediante el uso de \texttt{setARP}) que venían provistas en el archivo original. Debido a que el entorno y la resolución de direcciones ahora requieren de un descubrimiento puramente dinámico delegado a nuestro controlador SDN, las tablas ARP de los hosts inician vacías, permitiendo probar la lógica de intercepción y resolución ARP de manera realista.


% --- SECCIÓN 3 ---
\section{Manejo de ARP}
\label{sec:manejo_arp}

Para prescindir completamente de las configuraciones estáticas (como la instrucción \texttt{setARP} en \texttt{topo.py}), se implementó un mecanismo de resolución de direcciones ARP dinámico e inteligente en el controlador POX a través de una clase dedicada, \texttt{ArpHandler}.

Cuando un host intenta enviar un paquete y desconoce la dirección MAC de su puerta de enlace (la interfaz del router), envía un paquete ARP Request en \textit{broadcast}. El controlador intercepta estos paquetes e inspecciona el encabezado. Si el ARP Request está destinado a la IP de la interfaz privada del router (\texttt{192.168.1.254}) o a su interfaz pública (\texttt{200.0.0.254}), el controlador asume su rol de \textit{Default Gateway} y genera de inmediato un paquete ARP Reply informando la dirección MAC correspondiente (\texttt{PRIVATE\_MAC} o \texttt{PUBLIC\_MAC}). Esto permite al host emisor completar su tabla ARP local y comenzar a transmitir el tráfico IP.

Adicionalmente, el controlador aprende de forma proactiva observando todo el tráfico ARP entrante (tanto Requests como Replies). Cuando procesa estas tramas, extrae la información del emisor y registra o actualiza la relación entre la dirección IP y la dirección MAC en una tabla ARP dinámica interna (\texttt{arp\_table}).

El aspecto más crítico de este mecanismo surge cuando el switch debe enrutar un paquete IP hacia un destino (ya sea hacia el servidor público exterior o hacia un host privado interno) y desconoce su dirección MAC y su puerto físico. Para no bloquear el hilo de ejecución del controlador, se diseñó una resolución \textbf{asincrónica}:
\begin{enumerate}
    \item Se consulta la IP de destino en la tabla ARP del \texttt{ArpHandler} para obtener su MAC y puerto.
    \item Si la información no existe, se suspende momentáneamente el procesamiento de ese flujo. El paquete IP original (o evento) se almacena en una cola de espera (\texttt{waiting\_queue}), agrupado bajo la IP de destino que se desea resolver.
    \item El controlador genera y envía proactivamente su propio ARP Request. Si el tráfico se dirige a la red pública, el request sale por el puerto WAN conocido (\texttt{PUBLIC\_PORT}). Sin embargo, si el tráfico ingresa hacia la red privada y se desconoce la ubicación exacta del host, el request se envía mediante una inundación (\textit{flood} vía \texttt{of.OFPP\_FLOOD}) por todos los puertos locales para localizar físicamente al dispositivo.
    \item Dado que el procesamiento en POX está orientado a eventos, el controlador queda liberado para procesar otras tareas sin interrupciones.
    \item Cuando finalmente se recibe el paquete ARP Reply del dispositivo, el controlador aprende y registra tanto su MAC como su puerto de conexión en la tabla. Inmediatamente después, desencola todos los eventos que estaban en pausa esperando por dicha IP y los reinyecta en el flujo principal de procesamiento (\texttt{\_handle\_PacketIn}), despachándolos en orden estricto de llegada (FIFO).
\end{enumerate}

Este sistema garantiza no solo la eliminación de las dependencias estáticas en la topología, sino que asegura un enrutamiento totalmente dinámico donde no se pierden paquetes IP legítimos mientras se lleva a cabo el proceso de descubrimiento de red.

\subsection{Limpieza de Entradas ARP Obsoletas}
Para evitar que la tabla ARP interna del controlador acumule entradas de hosts que ya no están activos, se implementó un mecanismo de \textit{reap} en el método \texttt{\_reap\_arp\_table()}. Cada vez que se aprende o refresca una entrada (invocando \texttt{learn\_arp\_entry()}), se recorre la tabla y se eliminan todas aquellas entradas cuyo \textit{timestamp} supere los \texttt{ENTRY\_TIMEOUT} segundos (10\,s, mismo valor utilizado para el \texttt{idle\_timeout} de los flujos OpenFlow). De esta forma, el plano de control y el plano de datos se mantienen sincronizados: si el switch elimina un flujo por inactividad, la entrada ARP correspondiente también será purgada en la próxima operación de aprendizaje.


% --- SECCIÓN 4 ---
\section{Diseño de la Tabla NAT}
\label{sec:diseno_tabla_nat}

\subsection{Estructura de Datos e Identificación de Conexiones}
Para implementar el NAT con traducción de puertos (PAT) y soportar múltiples conexiones simultáneas, se diseñó la clase \texttt{NatHandler}, encargada de administrar el estado de las sesiones y la asignación dinámica de puertos. 

Dado que el controlador debe inspeccionar la capa de transporte (TCP/UDP) y traducir el tráfico en ambas direcciones de forma eficiente, se utilizaron dos diccionarios (tablas hash) en Python que funcionan como índices bidireccionales:

\begin{itemize}
    \item \textbf{Tabla Directa (\texttt{nat\_table}):} Se consulta cuando un paquete sale de la red privada hacia la pública. Su clave de búsqueda es una tupla que identifica unívocamente la sesión de origen: \texttt{(protocolo, IP\_privada, puerto\_privado)}. El valor almacenado contiene el \texttt{puerto\_público} asignado, la IP pública del router y un \textit{timestamp}.
    \item \textbf{Tabla Inversa (\texttt{reverse\_nat}):} Se consulta cuando el servidor público responde. La clave de búsqueda es la tupla \texttt{(protocolo, IP\_pública, puerto\_público\_destino)}. Al consultar este diccionario con el puerto de destino del paquete entrante, el router recupera en tiempo constante ($O(1)$) la IP y el puerto del host privado original para armar el flujo de retorno.
\end{itemize}

El \texttt{NatHandler} mantiene un conjunto (\textit{Set}) con los puertos públicos disponibles (desde el 1024). Al interceptar un flujo nuevo, se asigna el menor puerto disponible, asegurando que no existan colisiones entre las conexiones de diferentes hosts.

\subsection{Ejemplo Conceptual de la Tabla}
A continuación, se presenta un ejemplo de cómo se estructuran las asignaciones al establecerse dos conexiones en paralelo:

\begin{table}[h!]
\centering
\begin{tabular}{ccccc}
\toprule
\textbf{Protocolo} & \textbf{IP Privada} & \textbf{Puerto Privado} & \textbf{IP Pública} & \textbf{Puerto Público} \\
\midrule
TCP & 192.168.1.2 & 50001 & 200.0.0.254 & 1024 \\
UDP & 192.168.1.3 & 40002 & 200.0.0.254 & 1025 \\
\bottomrule
\end{tabular}
\caption{Ejemplo de asignación de puertos en las tablas bidireccionales del PAT.}
\label{tab:ejemplo_nat}
\end{table}

\subsection{Limpieza de Entradas NAT Expiradas}
\label{subsec:gc_nat}

Para evitar el agotamiento del pool de puertos públicos y mantener la coherencia con el plano de datos, se diseñó un mecanismo de \textit{garbage collection} dentro de \texttt{NatHandler}. El método \texttt{\_reap\_nat\_table()} itera sobre la \texttt{nat\_table} y elimina toda entrada cuyo \textit{timestamp} supere los \texttt{ENTRY\_TIMEOUT} segundos (10\,s). Por cada entrada eliminada se remueve también su contraparte en \texttt{reverse\_nat} y se libera el puerto público asociado mediante \texttt{release\_public\_port()}, devolviéndolo al conjunto de puertos disponibles para futuras conexiones.

El reap se ejecuta de manera \textit{triggerada} (no por temporizador) en los dos puntos de entrada a la tabla NAT:
\begin{itemize}
    \item \texttt{create\_pat\_entry()}: antes de crear o consultar una traducción saliente.
    \item \texttt{get\_reverse\_entry()}: antes de buscar una traducción entrante.
\end{itemize}

Esta estrategia cubre la mayoría de los casos prácticos, pero presenta una limitación conocida: si una conexión expira en el switch (por \texttt{idle\_timeout}) y posteriormente no se genera tráfico entrante ni saliente que pase por el controlador, la entrada NAT correspondiente nunca se reapará hasta que alguna de las dos funciones mencionadas sea invocada por una conexión completamente distinta. En escenarios de baja actividad, esto podría retrasar la liberación de puertos públicos. Una posible mejora futura sería incorporar un temporizador periódico independiente que ejecute el reap de forma asíncrona, desacoplando la limpieza de los puntos de consulta del mapa.


% --- SECCIÓN 5 ---
\section{Estrategia de Instalación de Flujos}
\label{sec:estrategia_flujos}

Cada vez que el controlador recibe el primer \textit{PacketIn} de una conexión (ya sea saliente o entrante), en lugar de traducir paquete por paquete, instala un par de flujos OpenFlow en el switch (\texttt{ofp\_flow\_mod}) para que los paquetes subsiguientes sean procesados directamente en el plano de datos, sin intervención del controlador. Esto minimiza el \textit{overhead} de comunicación entre el switch y el controlador, reduciendo drásticamente la latencia de los paquetes siguientes.

\subsection{Flujo Saliente (Privado → Público)}
\label{subsec:flow_out}

Cuando un paquete IP proveniente de la red privada llega al controlador, se instala un flujo que matchea los siguientes campos de la cabecera del paquete original:

\begin{verbatim}
  in_port, dl_type=0x800, nw_src, nw_dst, nw_proto, tp_src, tp_dst
\end{verbatim}

Las acciones aplicadas por el switch al matchear este flujo son:

\begin{enumerate}
    \item Reemplazar la IP de origen por la IP pública del router (\texttt{set\_nw\_src}).
    \item Reemplazar el puerto de origen por el puerto público asignado por PAT (\texttt{set\_tp\_src}).
    \item Reemplazar la MAC de origen por la MAC de la interfaz pública (\texttt{set\_dl\_src}).
    \item Reemplazar la MAC de destino por la MAC del \textit{gateway} público (\texttt{set\_dl\_dst}).
    \item Enviar el paquete por el puerto público (\texttt{output:PUBLIC\_PORT}).
\end{enumerate}

\subsection{Flujo Entrante (Público → Privado)}
\label{subsec:flow_in}

Junto con el flujo saliente se instala un segundo flujo para el camino de retorno, de modo que la respuesta del servidor público también sea traducida automáticamente por el switch. Este flujo matchea:

\begin{verbatim}
  in_port=PUBLIC_PORT, dl_type=0x800, nw_src=<IP_destino_original>,
  nw_dst=PUBLIC_IP, nw_proto, tp_src=<puerto_destino_original>,
  tp_dst=<puerto_público_asignado>
\end{verbatim}

Y ejecuta las acciones inversas:

\begin{enumerate}
    \item Reemplazar la IP de destino por la IP privada original (\texttt{set\_nw\_dst}).
    \item Reemplazar el puerto de destino por el puerto privado original (\texttt{set\_tp\_dst}).
    \item Reemplazar la MAC de origen por la MAC de la interfaz privada (\texttt{set\_dl\_src}).
    \item Reemplazar la MAC de destino por la MAC del host privado (\texttt{set\_dl\_dst}).
    \item Enviar el paquete por el puerto privado de entrada (\texttt{output:<in\_port>}).
\end{enumerate}

\subsection{Misma Estrategia para Tráfico Entrante}
\label{subsec:flow_in_reverse}

Cuando el primer paquete de una conexión entrante (respuesta del servidor público que coincide con una entrada PAT existente) llega al controlador, se aplica exactamente la misma lógica pero con los roles intercambiados: primero se instala el flujo que traduce de público a privado y luego el flujo complementario para que los próximos paquetes de la subsecuente comunicación saliente también sean traducidos en el switch.

\subsection{\textit{Idle Timeout} y Liberación de Recursos}
\label{subsec:timeout}

Todos los flujos se instalan con \texttt{idle\_timeout = ENTRY\_TIMEOUT} (10 segundos), sin \texttt{hard\_timeout}. Esto significa que un flujo permanece activo mientras haya tráfico en la conexión, y se elimina automáticamente tras 10 segundos de inactividad. La elección de un valor bajo permite:

\begin{itemize}
    \item Liberar rápidamente las entradas en la tabla de flujos del switch (un recurso finito y costoso).
    \item Sincronizar de forma natural la limpieza del plano de datos con la limpieza del plano de control: el mismo valor de \texttt{ENTRY\_TIMEOUT} se utiliza en los métodos \texttt{\_reap\_nat\_table()} y \texttt{\_reap\_arp\_table()} para purgar las entradas obsoletas en el controlador.
    \item Evitar que conexiones efímeras (como las de navegación web o consultas DNS) dejen residuos en las tablas del switch.
\end{itemize}

No se utiliza \texttt{hard\_timeout} porque eliminar un flujo de forma forzosa mientras la conexión aún está activa generaría \textit{PacketIn} innecesarios al controlador, degradando el rendimiento.


% --- SECCIÓN 6 ---
\section{Decisiones de Diseño}
\label{sec:decisiones_diseno}

\subsection{Separación de Responsabilidades en Handlers}
\label{subsec:handlers}

Se optó por modelar cada funcionalidad del plano de control en una clase independiente: \texttt{ArpHandler} (resolución de direcciones), \texttt{NatHandler} (traducción PAT para TCP/UDP) e \texttt{IcmpHandler} (traducción de IDs para ICMP). El orquestador \texttt{ProtoRouter} las instancia como atributos propios y las utiliza sin conocer los detalles internos de cada una. Esta decisión sigue el principio de responsabilidad única (SRP): cada clase encapsula su propio estado (tablas, pools de recursos) y su lógica asociada, lo que facilita la lectura, el mantenimiento y eventuales pruebas unitarias.

\subsection{Tabla NAT Bidireccional}
\label{subsec:tabla_bidireccional}

Para soportar la traducción en ambos sentidos con complejidad $O(1)$, se implementaron dos diccionarios paralelos:

\begin{itemize}
    \item \textbf{\texttt{nat\_table}}: clave \texttt{(protocolo, IP\_privada, puerto\_privado)} $\rightarrow$ datos de la traducción saliente.
    \item \textbf{\texttt{reverse\_nat}}: clave \texttt{(protocolo, IP\_pública, puerto\_público)} $\rightarrow$ datos de la traducción entrante.
\end{itemize}

Con un solo diccionario, la búsqueda inversa (necesaria cuando llega un paquete del servidor público) requeriría recorrer todas las entradas en $O(n)$. La duplicación controlada de índices es una compensación aceptable de memoria por velocidad de consulta, crítica para el procesamiento en tiempo real de paquetes.

\subsection{Algoritmo de Asignación de Puertos Públicos}
\label{subsec:asignacion_puertos}

El método \texttt{allocate\_public\_port()} selecciona el puerto disponible más bajo mediante \texttt{min()} sobre el conjunto \texttt{public\_port\_pool}. Se descartaron alternativas como:

\begin{itemize}
    \item \textbf{Asignación secuencial con índice}: requeriría mantener un contador global y lógica de \textit{wraparound}, además de gestionar huecos dejados por puertos liberados.
    \item \textbf{Asignación aleatoria}: introduce impredecibilidad, lo que dificulta la depuración.
\end{itemize}

La función \texttt{min()} sobre un \texttt{set} produce un comportamiento determinista: los puertos se reasignan desde el valor más bajo disponible, llenando naturalmente los huecos dejados por conexiones finalizadas. Esto replica de forma simplificada la política de puertos efímeros de los sistemas operativos.

\subsection{Pool de Puertos como \texttt{set} de Python}
\label{subsec:pool_set}

El conjunto de puertos públicos disponibles se implementa como un \texttt{set} de Python con el rango $[1024, 65535]$. Se privilegiaron los \texttt{set} sobre listas o colas porque:
\begin{itemize}
    \item Las operaciones de inserción (\texttt{add}) y eliminación (\texttt{remove}) son $O(1)$.
    \item La estructura garantiza ausencia de duplicados, incluso si \texttt{release\_public\_port()} se invocara múltiples veces sobre el mismo puerto por error.
\end{itemize}
El rango comienza en 1024 para excluir puertos privilegiados (0--1023) y reservados (por ejemplo, 22, 80, 443), evitando colisiones con servicios estándar.

\subsection{Estrategia de Limpieza de Entradas Expiradas}
\label{subsec:reap}

Se implementó un mecanismo de \textit{garbage collection} triggerado (no por temporizador) para depurar las tablas NAT, ARP e ICMP. Cada método público de consulta o creación ejecuta el reap antes de operar:

\begin{itemize}
    \item \texttt{create\_pat\_entry()} y \texttt{get\_reverse\_entry()} invocan \texttt{\_reap\_nat\_table()}.
    \item \texttt{learn\_arp\_entry()} invoca \texttt{\_reap\_arp\_table()}.
    \item Los métodos análogos de \texttt{IcmpHandler} siguen el mismo patrón.
\end{itemize}

Se descartó un temporizador periódico independiente para mantener el diseño simple y acoplado al flujo de eventos. La limitación conocida es que, en escenarios de inactividad total donde no se reciben \textit{PacketIn}, las entradas expiradas pueden permanecer hasta que una nueva conexión active el reap. Para un entorno académico con tráfico continuo durante las pruebas, esta estrategia es suficiente.

\subsection{Sincronización de Timeouts con \texttt{ENTRY\_TIMEOUT}}
\label{subsec:entry_timeout}

Se definió la constante \texttt{ENTRY\_TIMEOUT = 10} (segundos) que se utiliza tanto para el \texttt{idle\_timeout} de los flujos OpenFlow como para la expiración de las entradas en los tres handlers (NAT, ARP e ICMP). Esta decisión garantiza coherencia temporal entre el plano de datos (switch) y el plano de control (controlador): cuando el switch elimina un flujo por inactividad, el controlador purga el estado asociado aproximadamente en el mismo intervalo. No se emplea \texttt{hard\_timeout} para evitar que un flujo se elimine forzosamente mientras la conexión sigue activa, lo que generaría \textit{PacketIn} innecesarios.

\subsection{Resolución ARP Asincrónica con \texttt{waiting\_queue}}
\label{subsec:arp_async}

Cuando se necesita resolver una dirección MAC y la tabla ARP no contiene la entrada, el controlador no bloquea el procesamiento. En lugar de ello:
\begin{enumerate}
    \item Encola el evento (\texttt{PacketIn}) original en \texttt{waiting\_queue[IP\_destino]}.
    \item Envía proactivamente un ARP Request (por el puerto público conocido o mediante flooding en la red privada).
    \item Retorna el control al reactor de POX para procesar otros eventos.
    \item Al recibir el ARP Reply, aprende la entrada y desencola todos los eventos pendientes para reinyectarlos en el flujo principal.
\end{enumerate}
Este diseño es obligatorio en un framework orientado a eventos como POX, donde no es posible realizar esperas bloqueantes. Además, garantiza que ningún paquete IP legítimo se pierda durante la resolución.

\subsection{Instalación de Pares de Flujos Bidireccionales y Traducción del Primer Paquete}
\label{subsec:flujos_pares}

Por cada nueva conexión se instalan dos flujos OpenFlow (saliente y entrante) de forma simultánea. Esto minimiza la intervención del controlador: una vez instalados, todos los paquetes subsiguientes se traducen directamente en el switch. Como el primer paquete de la conexión no debe esperar a que los flujos se instalen para ser reenviado, el controlador lo traduce en software y lo envía inmediatamente mediante un \texttt{ofp\_packet\_out}. Los paquetes siguientes son manejados por los flujos en el plano de datos.

\subsection{Delegación del Cálculo de Checksums}
\label{subsec:checksums}

Después de modificar direcciones IP, puertos o IDs ICMP, se asignan los campos \texttt{csum} a \texttt{None}. POX recalcula automáticamente los checksums IP y de transporte al empaquetar la trama. Esta decisión evita implementar manualmente el cálculo de checksums y reduce la probabilidad de errores en la traducción.

\subsection{Manejo de ICMP como Protocolo Especial}
\label{subsec:icmp}

ICMP no posee puertos de transporte; utiliza un \texttt{icmp\_id} (campo \texttt{id} del \texttt{echo}) para identificar sesiones. Se creó un \texttt{IcmpHandler} independiente con su propio pool de IDs, análogo estructural al \texttt{NatHandler}. La función \texttt{\_get\_transport\_fields()} extrae el \texttt{icmp\_id} tanto como \texttt{src\_port} como \texttt{dst\_port}, lo que simplifica el manejo simétrico en los métodos \texttt{\_handle\_outgoing} y \texttt{\_handle\_incoming}.

\subsection{Filtrado de Protocolos No Soportados}
\label{subsec:filtrado}

Se definió la tupla \texttt{SUPPORTED\_PROTOCOLS = ("TCP", "UDP", "ICMP")}. Cualquier paquete IP con un protocolo de transporte diferente es descartado. Esta decisión se fundamenta en que PAT requiere inspeccionar puertos de transporte (TCP/UDP) o IDs (ICMP); protocolos como SCTP, GRE o IPIP carecen de este concepto y no pueden ser traducidos con este esquema.

\subsection{Refactorización con Métodos Auxiliares}
\label{subsec:refactor}

El código de \texttt{ProtoRouter} se organizó en métodos auxiliares reutilizables para evitar duplicación y mejorar la legibilidad:

\begin{itemize}
    \item \texttt{\_install\_flow()}: recibe un diccionario de campos match y una lista de acciones, construye y envía el \texttt{ofp\_flow\_mod}. Esto eliminó la repetición del bloque de instalación de flujos que antes aparecía cuatro veces.
    \item \texttt{\_resolve\_mac()} y \texttt{\_resolve\_mac\_and\_port()}: encapsulan la lógica de consulta ARP, encolado y envío de ARP Request.
    \item \texttt{\_clear\_checksums()}: centraliza la invalidación de checksums.
    \item \texttt{\_send\_packet()}: construye y envía el \texttt{ofp\_packet\_out} con el logging correspondiente.
\end{itemize}

\subsection{Uso de Constantes Descriptivas}
\label{subsec:constantes}

Todos los parámetros configurables (subred privada, IPs, MACs, puertos, timeouts) se definen como constantes en el ámbito del módulo, eliminando valores mágicos dispersos en el código. Esto facilita la modificación de la configuración sin necesidad de buscar literales en el cuerpo de las funciones.


% --- SECCIÓN 7 ---
\section{Problemas Encontrados y Soluciones Implementadas}
\label{sec:problemas_soluciones}

\begin{small}
\color{blue}
\textbf{Nota Provisoria para el Desarrollo:} 
Sección obligatoria para demostrar los dolores de cabeza del TP. Pueden listar problemas típicos como: el orden en el que se parsean los paquetes (ej. lidiar con paquetes que no eran IPv4 o ARP), problemas de concurrencia/colisión al asignar puertos rápidos, o comportamientos inesperados de los timeouts de OpenFlow, y cómo lograron resolverlos.
\end{small}

\vspace{0.5cm}
[Escribir aquí los problemas y soluciones...]


% --- SECCIÓN 8 ---
\section{Pruebas Realizadas}
\label{sec:pruebas}

\begin{small}
\color{blue}
\textbf{Nota Provisoria para el Desarrollo:} 
Acá vuelcan los resultados de la simulación. En las pruebas de demostración se suele exigir cambiar IPs/MACs para verificar que nada esté hardcodeado. Mencionen el uso de herramientas como Wireshark, Mininet CLI, iperf o tcpdump para validar el tráfico.
\end{small}

\subsection{Conectividad Básica}
[Detallar la prueba inicial de ping entre hosts para verificar el correcto funcionamiento del ARP dinámico y el NAT básico]

\subsection{Múltiples Clientes}
[Demostrar el comportamiento de la red privada expandida con más hosts (h2, h3, h4...) realizando consultas simultáneas hacia la red pública]

\subsection{Diferentes Tipos de Tráfico}
[Mostrar que la solución soporta exitosamente flujos TCP y UDP de forma aislada o combinada (ej. pruebas con iperf o wget/curl)]


% --- SECCIÓN 9 ---
\section{Conclusión}
\label{sec:conclusion}

\begin{small}
\color{blue}
\textbf{Nota Provisoria para el Desarrollo:} 
Cierre académico del informe. Hablen sobre los beneficios de SDN (flexibilidad para programar la lógica del NAT directo desde software) frente a las limitaciones observadas en el rendimiento por el overhead del primer Packet-In, consolidando los conceptos aprendidos sobre la separación del plano de control y de datos.
\end{small}

\vspace{0.5cm}
[Escribir aquí la conclusión del trabajo práctico...]

\end{document}